\section{Perception output}
\label{sec:tt_perception}

Each family of sensors is processed by a dedicated pipeline, running as an independent
process, and all the pipelines publish their output with the same structure, i.e.\ a list of
detections, each one representing a possible opponent vehicle and containing the acquisition
time of the sensor data, the source that produced it, the position of the object and, when
available, its heading and range rate. All the detections reach the tracker in the same
reference frame and, thanks to this common interface, the tracker treats all the sources in
the same way, the only differences being the measured quantities and their uncertainty.

All the pipelines exploit two pieces of prior information on the circuit. The first is a
binary map of the drivable area, which is used to discard the detections that fall outside the track boundaries. The second is the database of reference trajectories: when the sensor cannot measure the heading of a detection, it is taken from the closest
point of the closest trajectory.

\paragraph{LiDAR.} The point cloud of the LiDARs is processed by two complementary pipelines,
a geometric one, based on Clustering, and one based on a neural network, PointPillars.

In the Clustering pipeline, the points belonging to the ground are first removed, and only the
remaining points that fall inside the drivable area are kept. These points are grouped by a
density-based clustering algorithm (DBSCAN), and the clusters whose extent is not compatible
with the dimensions of a car are discarded. Since the LiDAR samples only the surfaces facing
the sensor, the points of each cluster depend on the visible part of the vehicle, and their
centroid is shifted towards the visible corner or side. Once this corner is identified, the
cluster is therefore shifted according to the nominal dimensions of the vehicle, so that its
position corresponds to the center of the car.

PointPillars, instead, is based on a deep neural network for three-dimensional object
detection, which operates on the same point cloud. The cloud is divided into vertical columns,
the pillars, whose points are encoded into features and processed by the network as an image.
For each detected object, the network returns an oriented three-dimensional bounding box,
i.e.\ its center, dimensions and heading, together with a confidence score. Since the network
learns the shape of the vehicle from the training data, it can estimate the center and the
heading of the vehicle also when it is only partially visible.

\paragraph{RaDAR.} At every cycle, each RaDAR returns a list of reflection points, each one
described by its range, azimuth, elevation and range rate. The points at ranges and heights
not compatible with a vehicle on track are discarded, and the remaining ones are grouped by a
spatial clustering whose parameters depend on the range, since the number of reflections
returned by a vehicle decreases with its distance and fewer points are therefore required to
form a cluster at longer range. The clusters are then filtered with the map of the drivable
area, and each of them generates a detection whose position is the centroid of its points and
whose range rate is their mean.

\paragraph{Camera.} The camera pipeline detects the vehicles in the images with a YOLO
neural network, which returns a two-dimensional bounding box for each of them. The position
of the vehicle is then reconstructed either from the camera alone, estimating the distance
from the apparent height of the bounding box, or by projecting
the LiDAR points onto the image and exploiting the points that fall inside the bounding box
to locate the vehicle more accurately.



% =====================================================================
\section{Prediction}
